\documentclass[11pt,a4paper]{article}
% =====================================================================
%  Gemeinsame Vorlage der Arbeitsblätter "Graphen" – Informatik 11 (NTG)
%  (gleiches Layout wie informatik/10/java/arbeitsblaetter/ab-vorlage.tex)
%  Einbinden in jedem Blatt direkt nach \documentclass: % =====================================================================
%  Gemeinsame Vorlage der Arbeitsblätter "Graphen" – Informatik 11 (NTG)
%  (gleiches Layout wie informatik/10/java/arbeitsblaetter/ab-vorlage.tex)
%  Einbinden in jedem Blatt direkt nach \documentclass: % =====================================================================
%  Gemeinsame Vorlage der Arbeitsblätter "Graphen" – Informatik 11 (NTG)
%  (gleiches Layout wie informatik/10/java/arbeitsblaetter/ab-vorlage.tex)
%  Einbinden in jedem Blatt direkt nach \documentclass: \input{ab-vorlage}
%  Übersetzen mit XeLaTeX, LuaLaTeX, pdfLaTeX oder Tectonic.
%
%  Blatt:   xelatex ab-3-1-grundbegriffe.tex
%  Lösung:  xelatex ab-3-1-grundbegriffe-lsg.tex
%           (Datei enthält nur: \def\mitloesung{}\input{ab-3-1-grundbegriffe})
%           (füllt alle Lücken rot aus; siehe \lsg, \lsglinien, \lsgoder)
% =====================================================================
\usepackage[a4paper,left=18mm,right=18mm,top=26mm,bottom=17mm,headheight=15mm,headsep=4mm,footskip=8mm]{geometry}
\usepackage{iftex}
\ifPDFTeX
  \usepackage[T1]{fontenc}
  \usepackage[utf8]{inputenc}
  \usepackage{helvet}
  \usepackage[scaled=0.86]{beramono}
  \renewcommand{\familydefault}{\sfdefault}
\else
  \usepackage{fontspec}
  \setmainfont{texgyreheros}[Extension=.otf,UprightFont=*-regular,BoldFont=*-bold,ItalicFont=*-italic,BoldItalicFont=*-bolditalic]
  \setmonofont{DejaVuSansMono}[Extension=.ttf,UprightFont=*,BoldFont=*-Bold,ItalicFont=*-Oblique,BoldItalicFont=*-BoldOblique,Scale=0.86]
\fi
\usepackage[ngerman,shorthands=off]{babel}
\usepackage{xcolor,graphicx,tabularx,array,enumitem,fancyhdr,tikz,listings,multicol,calc,etoolbox}
\usepackage[most]{tcolorbox}
\usetikzlibrary{arrows.meta,positioning,calc,decorations.pathreplacing,shapes.geometric}
\usepackage[hidelinks]{hyperref}
\setlength{\parindent}{0pt}
\setlength{\parskip}{3pt}
\setlist{nosep,leftmargin=*}

% ---------- Lösungsschalter ----------
\newif\ifloesung
\ifdefined\mitloesung\loesungtrue\fi
\newcommand{\lsgfarbe}{\color{red!75!black}}

% ---------- Kopf- und Fußzeile (einheitliche Form) ----------
\newcommand{\lehrkraft}{Hau}
\newcommand{\themenbereich}{3. Graphen}
\newcommand{\abnummer}{}
\pagestyle{fancy}\fancyhf{}
\renewcommand{\headrulewidth}{0.4pt}
\fancyhead[L]{\small Klasse 11\\Datum:}
\fancyhead[C]{\small Themenbereich:\\\themenbereich}
\fancyhead[R]{\small \lehrkraft\ifloesung\\{\lsgfarbe\bfseries Lösung}\fi}
\fancyfoot[C]{\scriptsize edu-mrh.de\,·\,Informatik 11\,·\,Blatt \abnummer\ifloesung\ (Lösung)\fi\,·\,Seite \thepage}
\newcommand{\abtitel}[2]{\renewcommand{\abnummer}{#1}\begin{center}\Large\bfseries #1\quad\underline{#2}\end{center}\vspace{-1mm}}

% ---------- Boxen ----------
\newenvironment{aufgabe}[2][]{\begin{tcolorbox}[enhanced,breakable,colback=white,colframe=black,boxrule=0.7pt,arc=3mm,left=2mm,right=2mm,top=1.5mm,bottom=1.5mm,before skip=2.5mm,after skip=2.5mm]\textbf{Aufgabe #2}\ifstrempty{#1}{}{ \textit{#1}}\par\vspace{0.6mm}}{\end{tcolorbox}}
\newtcolorbox{merke}[1][Merke]{enhanced,breakable,colback=green!5,colframe=green!40!black,boxrule=0.8pt,arc=2mm,left=2mm,right=2mm,top=1.5mm,bottom=1.5mm,before skip=2.5mm,after skip=2.5mm,title={#1},fonttitle=\bfseries,coltitle=white}
\newtcolorbox{hinweis}{enhanced,breakable,colback=yellow!12,colframe=orange!70!black,boxrule=0.6pt,arc=2mm,left=2mm,right=2mm,top=1mm,bottom=1mm,before skip=2mm,after skip=2mm}
\newtcolorbox{ausblick}[1][Ausblick]{enhanced,breakable,colback=teal!5,colframe=teal!60!black,boxrule=0.8pt,arc=2mm,left=2mm,right=2mm,top=1.5mm,bottom=1.5mm,before skip=2.5mm,after skip=2.5mm,title={#1},fonttitle=\bfseries,coltitle=white}
\newcommand{\web}[1]{{\small\color{gray}\textit{Website: #1}}}

% ---------- Lücken, Linien, Karos ----------
\newcommand{\luecke}[1][3.5cm]{\rule[-0.5ex]{#1}{0.4pt}}
\newcommand{\linien}[1]{\par\vspace{1.5mm}\foreach \i in {1,...,#1}{\noindent\rule[0pt]{\linewidth}{0.3pt}\par\vspace{5.5mm}}\vspace{-3mm}}
\newcommand{\kariert}[1]{\par\vspace{1.5mm}\noindent\begin{tikzpicture}\draw[step=5mm,gray!55,thin] (0,0) grid (\linewidth-1pt,#1*5mm);\end{tikzpicture}\par}
\newcommand{\karobox}[2]{\begin{tikzpicture}\draw[step=5mm,gray!55,thin] (0,0) grid (#1,#2);\end{tikzpicture}}
\newcommand{\code}[1]{\texttt{#1}}

% Lücke mit Lösung: \lsg[Mindestbreite]{Lösungstext}
\newlength{\lsgbreite}
\newcommand{\lsg}[2][3.5cm]{\ifloesung\settowidth{\lsgbreite}{\,#2\,}\ifdim\lsgbreite<#1\setlength{\lsgbreite}{#1}\fi\underline{\makebox[\lsgbreite][l]{\lsgfarbe\,#2}}\else\luecke[#1]\fi}
% Schreiblinien mit Lösung: \lsglinien{Anzahl}{Lösungstext} (gleiche Höhe wie \linien)
\newcommand{\lsglinien}[2]{\ifloesung\par\vspace{1.5mm}\noindent\begin{minipage}[t][#1\dimexpr 6.9mm\relax][t]{\linewidth}\lsgfarbe #2\end{minipage}\par\vspace{-1.5mm}\else\linien{#1}\fi}
% Beliebiger Inhalt: \lsgoder{nur in der Lösung}{nur im Blatt}
\newcommand{\lsgoder}[2]{\ifloesung\leavevmode#1\else#2\fi}

% ---------- Ergänzung (aus informatik/10/java/arbeitsblaetter/ab-vorlage.tex) ----------
% Bestehende Makros oben bleiben unverändert, damit ältere Blätter gleich aussehen.
%  \restzeile / \lsgrest{Text}   Lücke bis zum Zeilenende
%  \lsgtext{Text}                nur in der Lösung sichtbar (z. B. Tabellenzellen)
%  \lueckentext                  größerer Zeilenabstand, Flattersatz für Absätze mit Lücken
%  \schreibzeile                 Strut für Tabellenzellen zum Ausfüllen (8 mm hoch)
%  \alt{9}, \neu{7}              überschriebener alter Wert / neuer Wert (Dijkstra-Tabellen)
\usepackage{colortbl,xspace}
\usepackage[normalem]{ulem}
\newcommand{\lsgrest}[1]{\leavevmode\unskip\hspace{1.5mm}\ifloesung\rlap{\raisebox{0.8pt}{\lsgfarbe\,#1}}\fi\leaders\hbox{\rule[-0.6ex]{0.5pt}{0.4pt}}\hfill\null\par}
\newcommand{\restzeile}{\lsgrest{}}
\newcommand{\lsgtext}[1]{\ifloesung\leavevmode{\lsgfarbe#1}\fi}
\newcommand{\lueckentext}{\linespread{1.5}\selectfont\raggedright}
\newcommand{\schreibzeile}{\rule[-2.8mm]{0pt}{8mm}}
\newcommand{\alt}[1]{\leavevmode{\color{blue!60!black}\sout{#1}}}
\newcommand{\neu}[1]{\leavevmode{\color{green!50!black}\bfseries #1}}
\colorlet{djaktuell}{orange!40}
\colorlet{djfertig}{blue!22}

% ---------- Graphen zeichnen ----------
\tikzset{
  knoten/.style={draw,rounded corners=3pt,fill=orange!12,minimum height=6mm,inner sep=3pt,font=\small},
  kknoten/.style={draw,circle,fill=orange!12,minimum size=7mm,inner sep=1pt,font=\small},
  kante/.style={thick},
  gkante/.style={thick,-{Stealth[length=2.6mm]}},
  dkante/.style={thick,{Stealth[length=2.6mm]}-{Stealth[length=2.6mm]}},
  gewicht/.style={font=\scriptsize,fill=white,inner sep=1pt,text=teal!60!black},
  lsgkante/.style={thick,red!75!black},
  lsggkante/.style={thick,red!75!black,-{Stealth[length=2.6mm]}},
  lsgknoten/.style={knoten,draw=red!75!black,text=red!75!black,fill=red!4},
}
% Adjazenzmatrix zum Ausfüllen: \adjmx[Lösungseinträge]{Anzahl}{Knotenliste}{Zellbreite}
% Lösungseinträge als Zeile/Spalte/Text, z. B. \adjmx[1/2/X,2/1/X]{3}{A,B,C}{7mm}
\newcommand{\adjmx}[4][]{\begin{tikzpicture}[x=#4,y=#4]
  \foreach \n [count=\i] in {#3} {\node[font=\small\bfseries] at (\i,0) {\n}; \node[font=\small\bfseries,anchor=east] at (0.45,-\i) {\n};}
  \foreach \k in {0,...,#2} {\draw[gray] (\k+0.5,-0.5) -- (\k+0.5,-#2-0.5); \draw[gray] (0.5,-\k-0.5) -- (#2+0.5,-\k-0.5);}
  \ifloesung\ifstrempty{#1}{}{\foreach \z/\s/\t in {#1} {\node[text=red!75!black,font=\small] at (\s,-\z) {\t};}}\fi
\end{tikzpicture}}

% ---------- Java-Listings ----------
\lstset{language=Java,basicstyle=\ttfamily\small,keywordstyle=\bfseries,commentstyle=\color{gray!80!black},stringstyle=\color{black},showstringspaces=false,columns=flexible,keepspaces=true,tabsize=3,frame=single,framerule=0.4pt,rulecolor=\color{black},xleftmargin=2mm,framexleftmargin=1mm,aboveskip=2mm,belowskip=2mm,breaklines=true,escapechar=§,
 literate={ä}{{\"a}}1 {ö}{{\"o}}1 {ü}{{\"u}}1 {Ä}{{\"A}}1 {Ö}{{\"O}}1 {Ü}{{\"U}}1 {ß}{{\ss}}1 {–}{{--}}1 {„}{{\glqq}}1 {“}{{\grqq}}1}
\lstnewenvironment{java}[1][]{\lstset{#1}}{}

%  Übersetzen mit XeLaTeX, LuaLaTeX, pdfLaTeX oder Tectonic.
%
%  Blatt:   xelatex ab-3-1-grundbegriffe.tex
%  Lösung:  xelatex ab-3-1-grundbegriffe-lsg.tex
%           (Datei enthält nur: \def\mitloesung{}\documentclass[11pt,a4paper]{article}
\input{ab-vorlage}
% Dörfer-Graph (Merkkasten Begriffe): Koordinaten wie im Begriffe-Explorer der Website
\newcommand{\doerferknoten}{%
  \node[knoten,font=\tiny,minimum height=4mm,inner sep=1.5pt] (A) at (4.4,3.3) {Altdorf};
  \node[knoten,font=\tiny,minimum height=4mm,inner sep=1.5pt] (W) at (1.2,2.9) {Weiler};
  \node[knoten,font=\tiny,minimum height=4mm,inner sep=1.5pt] (F) at (5.1,2.0) {Fischbach};
  \node[knoten,font=\tiny,minimum height=4mm,inner sep=1.5pt] (Z) at (2.9,2.0) {Ziegelstein};
  \node[knoten,font=\tiny,minimum height=4mm,inner sep=1.5pt] (N) at (0.9,0.95) {Neustadt};
  \node[knoten,font=\tiny,minimum height=4mm,inner sep=1.5pt] (B) at (4.4,0.8) {Burg};
  \node[knoten,font=\tiny,minimum height=4mm,inner sep=1.5pt] (R) at (2.6,0.2) {Rain};}
\begin{document}
\abtitel{3.1}{Graphen -- Grundbegriffe}

Hänsel und Gretel haben die Hexe überlistet. Doch die Brotkrumen, mit denen sie ihren Weg markiert hatten, haben die Vögel gefressen. Findest du trotzdem nach Hause?

\begin{aufgabe}[Hänsel und Gretel]{1}
Starte das Spiel auf der Website (Station 3.1).
\textbf{a)} Zeichne beim Spielen eine Karte des Waldes: jeder Ort ein Kästchen, jeder Weg eine Linie. Führt ein Weg nur in eine Richtung, zeichne einen Pfeil. Das Hexenhaus ist schon eingezeichnet.
\end{aufgabe}
\noindent\begin{tikzpicture}
\draw[step=5mm,gray!55,thin] (0,0) grid (17.4,7);
\node[knoten,fill=green!15] (hex) at (4,2.3) {Hexenhaus (Start)};
\ifloesung
  \node[lsgknoten] (warn) at (5.5,6.5) {Warnung};
  \node[lsgknoten] (wolf) at (10.5,6.5) {Werwolf};
  \node[lsgknoten] (stadt) at (1.6,5.2) {Stadt};
  \node[lsgknoten] (abz) at (5.5,5.2) {Abzweigung};
  \node[lsgknoten] (spinne) at (13,5.2) {Spinne};
  \node[lsgknoten] (baeck) at (1.6,3.6) {Bäckerei};
  \node[lsgknoten] (pilze) at (4,3.9) {Pilze};
  \node[lsgknoten] (w3) at (13,3.4) {Wald3};
  \node[lsgknoten] (w1) at (8,1.0) {Wald1};
  \node[lsgknoten] (stein) at (12.2,2.0) {Steinkreis};
  \node[lsgknoten] (w2) at (15,0.8) {Wald2};
  \node[lsgknoten] (heim) at (16.2,2.6) {Zuhause};
  \draw[lsgkante] (stadt) -- (baeck);
  \draw[lsgkante] (stadt) -- (pilze);
  \draw[lsgkante] (pilze) -- (abz);
  \draw[lsgkante] (pilze) -- (hex);
  \draw[lsggkante] (pilze) to[bend left=50] node[right,font=\scriptsize,text=red!75!black] {Pilze essen} (hex);
  \draw[lsgkante] (abz) -- (warn);
  \draw[lsggkante] (warn) -- (wolf);
  \draw[lsgkante] (abz) -- (spinne);
  \draw[lsgkante] (spinne) -- (w3);
  \draw[lsgkante] (hex) -- (w1);
  \draw[lsgkante] (w1) -- (w3);
  \draw[lsgkante] (w1) -- (w2);
  \draw[lsgkante] (w1) -- (stein);
  \draw[lsgkante] (w2) -- (w3);
  \draw[lsgkante] (w2) -- (stein);
  \draw[lsgkante] (w3) -- (stein);
  \draw[lsggkante] (w2) -- (heim);
  \node[text=red!75!black,font=\scriptsize,align=left,anchor=north west] at (0.2,1.3) {13 Knoten, 17 Kanten\\(14 ungerichtet, 3 gerichtet)};
\fi
\end{tikzpicture}

\textbf{b)} Am Ende zeigt dir das Spiel deinen Weg. Notiere ihn und die Anzahl der Schritte:
\lsglinien{1}{individuell; kürzester Weg: Hexenhaus -- Wald1 -- Wald2 -- Zuhause (3 Schritte)}

\begin{minipage}[t]{0.6\linewidth}
\begin{merke}[Graph und Kurzschreibweise $G = (V, E)$]\raggedright
Ein Graph besteht aus \lsg[2.3cm]{Knoten} (engl. \textit{vertices}) und \lsg[2.3cm]{Kanten} (engl. \textit{edges}). Eine Kante verbindet zwei Knoten. Im Spiel sind die Orte die \lsg[2cm]{Knoten}, die Wege die \lsg[2cm]{Kanten}.\\[1mm]
Kurzschreibweise $G = (V, E)$: $V$ ist die Menge der Knoten, $E$ die Menge der Kanten.
\begin{itemize}
\item $(A, B)$: \lsg[2.5cm]{gerichtete} Kante, nur von $A$ nach $B$
\item $\{A, B\}$: \lsg[2.5cm]{ungerichtete} Kante, steht für $(A, B)$ und \lsg[1.4cm]{$(B, A)$}
\end{itemize}
\end{merke}
\end{minipage}\hfill
\begin{minipage}[t]{0.37\linewidth}
\vspace{1mm}
\centering\small \textbf{Beispiel: Wer folgt wem?}\\[1mm]
\begin{tikzpicture}[x=0.95cm,y=0.9cm]
\node[knoten] (leon) at (0.7,2.2) {Leon};
\node[knoten] (letifa) at (3.3,2.4) {Letifa};
\node[knoten] (luise) at (4.3,1.1) {Luise};
\node[knoten] (hamza) at (2.4,0.3) {Hamza};
\node[knoten] (sabine) at (0.5,0.6) {Sabine};
\draw[dkante] (leon) -- (letifa);
\draw[gkante] (letifa) -- (luise);
\draw[gkante] (leon) -- (luise);
\draw[dkante] (leon) -- (hamza);
\draw[gkante] (hamza) -- (luise);
\end{tikzpicture}
\end{minipage}\\[1mm]
{\small $V = \{\text{Hamza}, \text{Leon}, \text{Letifa}, \text{Luise}, \text{Sabine}\}$\\
$E = \{\{\text{Leon}, \text{Letifa}\}, \{\text{Leon}, \text{Hamza}\}, (\text{Letifa}, \text{Luise}), (\text{Leon}, \text{Luise}), (\text{Hamza}, \text{Luise})\}$}

\begin{aufgabe}[Hänsel und Gretel]{1}
\textbf{c)} Schreibe den Ausschnitt deiner Karte mit den Knoten Hexenhaus, Pilze, Stadt, Bäckerei und Wald1 in Kurzschreibweise auf (5 Kanten).\\[2mm]
$V = \{$\lsg[15.4cm]{Hexenhaus, Pilze, Stadt, Bäckerei, Wald1}$\}$\\[3mm]
$E = \{$\lsg[15.4cm]{\{Hexenhaus, Pilze\}, (Pilze, Hexenhaus), \{Pilze, Stadt\}, \{Stadt, Bäckerei\}, \{Hexenhaus, Wald1\}}$\}$
\end{aufgabe}

\newpage
\begin{minipage}[t]{0.71\linewidth}
\begin{merke}[Pfade, Zyklen, Richtung, Gewicht]\raggedright\small
\textbf{Pfad (Weg):} Folge von Knoten, in der aufeinanderfolgende Knoten durch eine \lsg[1.8cm]{Kante} verbunden sind; keine Kante wird doppelt benutzt.\\[0.6mm]
\textbf{Einfacher Pfad:} Pfad, in dem sich kein \lsg[1.8cm]{Knoten} wiederholt.\\
Beispiel in A: \lsg[5.6cm]{Altdorf -- Weiler -- Ziegelstein -- Neustadt}\\[0.6mm]
\textbf{Länge eines Pfades:} ungewichtet die \lsg[3.2cm]{Anzahl der Kanten}, gewichtet die \lsg[3.2cm]{Summe der Gewichte}.\\[0.6mm]
\textbf{Zyklus:} einfacher Pfad, dessen \lsg[5cm]{Start- und Endknoten gleich sind}.\\
Beispiel in A: \lsg[5.6cm]{Weiler -- Ziegelstein -- Fischbach -- Weiler}\\
Graph mit Zyklus: \emph{zyklisch}, sonst \lsg[3.4cm]{azyklisch (zyklenfrei)}.\\[0.6mm]
\textbf{Gerichtet:} Kanten nur \lsg[3cm]{in Pfeilrichtung} (Pfeil).\ \textbf{Ungerichtet:} \lsg[2.8cm]{in beide Richtungen} (Strich).\\[0.6mm]
\textbf{Stark zusammenhängend:} Von \lsg[2.4cm]{jedem Knoten} gibt es einen Pfad zu \lsg[3.2cm]{jedem anderen Knoten}.\\
\textbf{Schwach zusammenhängend:} \lsg[5.4cm]{nur ohne Beachtung der Richtungen}\\ \lsg[\linewidth]{hängt alles zusammen. B ist schwach, denn von Rain führt keine Kante weg.}\\[0.6mm]
\textbf{Gewichtet:} Jede Kante trägt eine \lsg[2.6cm]{Zahl (Gewicht)}, z.\,B. \lsg[3.6cm]{Entfernung, Fahrzeit}.
\end{merke}
\end{minipage}\hfill
\begin{minipage}[t]{0.27\linewidth}
\vspace{1mm}
\textbf{\small A} {\scriptsize ungerichtet}\\
\begin{tikzpicture}[x=0.8cm,y=0.85cm]
\doerferknoten
\draw[kante] (A)--(W); \draw[kante] (F)--(A); \draw[kante] (W)--(F); \draw[kante] (W)--(Z);
\draw[kante] (Z)--(F); \draw[kante] (Z)--(N); \draw[kante] (B)--(Z); \draw[kante] (F)--(B);
\draw[kante] (N)--(R); \draw[kante] (B)--(R);
\end{tikzpicture}\\[2mm]
\textbf{\small B} {\scriptsize gerichtet}\\
\begin{tikzpicture}[x=0.8cm,y=0.85cm]
\doerferknoten
\draw[gkante] (A)--(W); \draw[gkante] (F)--(A); \draw[gkante] (W)--(F); \draw[gkante] (W)--(Z);
\draw[gkante] (Z)--(F); \draw[gkante] (Z)--(N); \draw[gkante] (B)--(Z); \draw[gkante] (F)--(B);
\draw[gkante] (N)--(R); \draw[gkante] (B)--(R);
\end{tikzpicture}
\end{minipage}

\begin{aufgabe}{2}
{\raggedright Zeichne $G = (V, E)$ mit $V = \{A, B, C, D, E, F\}$ und\\ $E = \{(A, B), (A, C), (A, E), (B, C), (C, C), (C, D), (E, C), (F, D)\}$. Beantworte die Fragen.\par}\vspace{1mm}
\begin{minipage}[t]{0.42\linewidth}
\vspace{0pt}
\begin{tikzpicture}
\draw[step=5mm,gray!55,thin] (0,0) grid (7,3);
\ifloesung
  \foreach \n/\x/\y/\t in {a/0.6/1.5/A, b/2/2.55/B, e/2/0.45/E, c/3.4/1.5/C, d/4.9/1.5/D, f/6.4/1.5/F} {
    \node[kknoten,minimum size=6mm,draw=red!75!black,text=red!75!black] (\n) at (\x,\y) {\t};}
  \draw[lsggkante] (a)--(b); \draw[lsggkante] (a)--(c); \draw[lsggkante] (a)--(e);
  \draw[lsggkante] (b)--(c); \draw[lsggkante] (c)--(d); \draw[lsggkante] (e)--(c);
  \draw[lsggkante] (f)--(d);
  \draw[lsggkante] (c) to[out=60,in=120,looseness=6] (c);
\fi
\end{tikzpicture}
\end{minipage}\hfill
\begin{minipage}[t]{0.55\linewidth}
\vspace{0pt}\small\raggedright
Gerichtet? Woran erkennbar? \lsg[3.9cm]{ja, Paare ( ) statt \{ \}}\\[1.5mm]
Ein Zyklus: \lsg[5.7cm]{(C, C), eine Schlinge (Länge 1)}\\[1.5mm]
Stark/schwach zusammenhängend? \lsg[3.6cm]{schwach: von D}\\[1.5mm]
\lsg[0.98\linewidth]{kommt man nirgends hin.}\\[1.5mm]
Länge des längsten einfachen Pfades: \lsg[3.2cm]{3 (A--B--C--D)}
\end{minipage}
\end{aufgabe}

\begin{aufgabe}[Zurück in den Wald]{3}
Deine Karte aus Aufgabe 1: \textbf{a)} gerichtet?\ \textbf{b)} ein Zyklus?\ \textbf{c)} stark oder schwach zusammenhängend?\ \textbf{d)} „Der Weg zurück … scheint verschwunden“: Ist die Kante Hexenhaus -- Wald1 gerichtet? Prüfe im Spiel.
\lsglinien{2}{\small a) gemischt; gerichtet nur (Pilze, Hexenhaus), (Warnung, Werwolf), (Wald2, Zuhause)\quad b) Wald1 -- Wald2 -- Wald3 -- Wald1\\ c) schwach: Von Zuhause/Werwolf führt keine Kante weg.\quad d) nein: Von Wald2 aus führt ein Weg zurück.}
\end{aufgabe}

\begin{aufgabe}[Ordne die Graphen ein.]{4}
\vspace{-1mm}
\small\renewcommand{\arraystretch}{1.15}
\begin{tabularx}{\linewidth}{|>{\centering\arraybackslash}m{5.6cm}|>{\centering\arraybackslash}m{1.7cm}|>{\centering\arraybackslash}m{1.1cm}|>{\centering\arraybackslash}m{1.1cm}|>{\centering\arraybackslash}m{1.1cm}|>{\centering\arraybackslash}X|}\hline
\textbf{Graph} & \scriptsize\textbf{zusammen-\newline hängend?}\newline stark/schwach/nein & \scriptsize\textbf{ge-\newline richtet?} & \scriptsize\textbf{ge-\newline wichtet?} & \scriptsize\textbf{zyklisch?} & \scriptsize\textbf{Länge eines längsten einfachen Pfades}\\\hline
\rule{0pt}{8mm}\textbf{I}\ \begin{tikzpicture}[baseline=(k5.base),x=0.85cm,y=0.55cm]
\foreach \n/\x/\y in {1/0/1,2/1/1,3/2/1,4/3/1,5/1/0,6/3/0} {\node[kknoten,minimum size=5mm,inner sep=0pt,font=\tiny] (k\n) at (\x,\y) {K\n};}
\draw[kante] (k1)--(k2)--(k3)--(k4); \draw[kante] (k2)--(k5);
\end{tikzpicture}
& \lsgoder{\lsgfarbe nein}{} & \lsgoder{\lsgfarbe nein}{} & \lsgoder{\lsgfarbe nein}{} & \lsgoder{\lsgfarbe nein}{} & \lsgoder{\lsgfarbe 3: K1--K2--K3--K4}{}\\\hline
\rule{0pt}{5mm}\textbf{II}\ $V = \{A, B, C\}$\newline $E =\{(A, B), (B, A), (C, B), (C, C)\}$
& \lsgoder{\lsgfarbe schwach}{} & \lsgoder{\lsgfarbe ja}{} & \lsgoder{\lsgfarbe nein}{} & \lsgoder{\lsgfarbe ja}{} & \lsgoder{\lsgfarbe 2: C$\to$B$\to$A}{}\\\hline
\rule{0pt}{8mm}\textbf{III}\ \begin{tikzpicture}[baseline=(k3.base),x=0.9cm,y=0.6cm]
\foreach \n/\x/\y in {1/0/1,2/2/1,3/1/0,4/3/0} {\node[kknoten,minimum size=5mm,inner sep=0pt,font=\tiny] (k\n) at (\x,\y) {K\n};}
\draw[kante] (k1) -- node[gewicht] {4} (k2); \draw[kante] (k2) -- node[gewicht] {1} (k3);
\draw[kante] (k3) -- node[gewicht] {5} (k1); \draw[kante] (k3) -- node[gewicht] {2} (k4);
\end{tikzpicture}
& \lsgoder{\lsgfarbe stark}{} & \lsgoder{\lsgfarbe nein}{} & \lsgoder{\lsgfarbe ja}{} & \lsgoder{\lsgfarbe ja}{} & \lsgoder{\lsgfarbe 11: K2--K1--K3--K4}{}\\\hline
\end{tabularx}
\end{aufgabe}
\end{document}
)
%           (füllt alle Lücken rot aus; siehe \lsg, \lsglinien, \lsgoder)
% =====================================================================
\usepackage[a4paper,left=18mm,right=18mm,top=26mm,bottom=17mm,headheight=15mm,headsep=4mm,footskip=8mm]{geometry}
\usepackage{iftex}
\ifPDFTeX
  \usepackage[T1]{fontenc}
  \usepackage[utf8]{inputenc}
  \usepackage{helvet}
  \usepackage[scaled=0.86]{beramono}
  \renewcommand{\familydefault}{\sfdefault}
\else
  \usepackage{fontspec}
  \setmainfont{texgyreheros}[Extension=.otf,UprightFont=*-regular,BoldFont=*-bold,ItalicFont=*-italic,BoldItalicFont=*-bolditalic]
  \setmonofont{DejaVuSansMono}[Extension=.ttf,UprightFont=*,BoldFont=*-Bold,ItalicFont=*-Oblique,BoldItalicFont=*-BoldOblique,Scale=0.86]
\fi
\usepackage[ngerman,shorthands=off]{babel}
\usepackage{xcolor,graphicx,tabularx,array,enumitem,fancyhdr,tikz,listings,multicol,calc,etoolbox}
\usepackage[most]{tcolorbox}
\usetikzlibrary{arrows.meta,positioning,calc,decorations.pathreplacing,shapes.geometric}
\usepackage[hidelinks]{hyperref}
\setlength{\parindent}{0pt}
\setlength{\parskip}{3pt}
\setlist{nosep,leftmargin=*}

% ---------- Lösungsschalter ----------
\newif\ifloesung
\ifdefined\mitloesung\loesungtrue\fi
\newcommand{\lsgfarbe}{\color{red!75!black}}

% ---------- Kopf- und Fußzeile (einheitliche Form) ----------
\newcommand{\lehrkraft}{Hau}
\newcommand{\themenbereich}{3. Graphen}
\newcommand{\abnummer}{}
\pagestyle{fancy}\fancyhf{}
\renewcommand{\headrulewidth}{0.4pt}
\fancyhead[L]{\small Klasse 11\\Datum:}
\fancyhead[C]{\small Themenbereich:\\\themenbereich}
\fancyhead[R]{\small \lehrkraft\ifloesung\\{\lsgfarbe\bfseries Lösung}\fi}
\fancyfoot[C]{\scriptsize edu-mrh.de\,·\,Informatik 11\,·\,Blatt \abnummer\ifloesung\ (Lösung)\fi\,·\,Seite \thepage}
\newcommand{\abtitel}[2]{\renewcommand{\abnummer}{#1}\begin{center}\Large\bfseries #1\quad\underline{#2}\end{center}\vspace{-1mm}}

% ---------- Boxen ----------
\newenvironment{aufgabe}[2][]{\begin{tcolorbox}[enhanced,breakable,colback=white,colframe=black,boxrule=0.7pt,arc=3mm,left=2mm,right=2mm,top=1.5mm,bottom=1.5mm,before skip=2.5mm,after skip=2.5mm]\textbf{Aufgabe #2}\ifstrempty{#1}{}{ \textit{#1}}\par\vspace{0.6mm}}{\end{tcolorbox}}
\newtcolorbox{merke}[1][Merke]{enhanced,breakable,colback=green!5,colframe=green!40!black,boxrule=0.8pt,arc=2mm,left=2mm,right=2mm,top=1.5mm,bottom=1.5mm,before skip=2.5mm,after skip=2.5mm,title={#1},fonttitle=\bfseries,coltitle=white}
\newtcolorbox{hinweis}{enhanced,breakable,colback=yellow!12,colframe=orange!70!black,boxrule=0.6pt,arc=2mm,left=2mm,right=2mm,top=1mm,bottom=1mm,before skip=2mm,after skip=2mm}
\newtcolorbox{ausblick}[1][Ausblick]{enhanced,breakable,colback=teal!5,colframe=teal!60!black,boxrule=0.8pt,arc=2mm,left=2mm,right=2mm,top=1.5mm,bottom=1.5mm,before skip=2.5mm,after skip=2.5mm,title={#1},fonttitle=\bfseries,coltitle=white}
\newcommand{\web}[1]{{\small\color{gray}\textit{Website: #1}}}

% ---------- Lücken, Linien, Karos ----------
\newcommand{\luecke}[1][3.5cm]{\rule[-0.5ex]{#1}{0.4pt}}
\newcommand{\linien}[1]{\par\vspace{1.5mm}\foreach \i in {1,...,#1}{\noindent\rule[0pt]{\linewidth}{0.3pt}\par\vspace{5.5mm}}\vspace{-3mm}}
\newcommand{\kariert}[1]{\par\vspace{1.5mm}\noindent\begin{tikzpicture}\draw[step=5mm,gray!55,thin] (0,0) grid (\linewidth-1pt,#1*5mm);\end{tikzpicture}\par}
\newcommand{\karobox}[2]{\begin{tikzpicture}\draw[step=5mm,gray!55,thin] (0,0) grid (#1,#2);\end{tikzpicture}}
\newcommand{\code}[1]{\texttt{#1}}

% Lücke mit Lösung: \lsg[Mindestbreite]{Lösungstext}
\newlength{\lsgbreite}
\newcommand{\lsg}[2][3.5cm]{\ifloesung\settowidth{\lsgbreite}{\,#2\,}\ifdim\lsgbreite<#1\setlength{\lsgbreite}{#1}\fi\underline{\makebox[\lsgbreite][l]{\lsgfarbe\,#2}}\else\luecke[#1]\fi}
% Schreiblinien mit Lösung: \lsglinien{Anzahl}{Lösungstext} (gleiche Höhe wie \linien)
\newcommand{\lsglinien}[2]{\ifloesung\par\vspace{1.5mm}\noindent\begin{minipage}[t][#1\dimexpr 6.9mm\relax][t]{\linewidth}\lsgfarbe #2\end{minipage}\par\vspace{-1.5mm}\else\linien{#1}\fi}
% Beliebiger Inhalt: \lsgoder{nur in der Lösung}{nur im Blatt}
\newcommand{\lsgoder}[2]{\ifloesung\leavevmode#1\else#2\fi}

% ---------- Ergänzung (aus informatik/10/java/arbeitsblaetter/ab-vorlage.tex) ----------
% Bestehende Makros oben bleiben unverändert, damit ältere Blätter gleich aussehen.
%  \restzeile / \lsgrest{Text}   Lücke bis zum Zeilenende
%  \lsgtext{Text}                nur in der Lösung sichtbar (z. B. Tabellenzellen)
%  \lueckentext                  größerer Zeilenabstand, Flattersatz für Absätze mit Lücken
%  \schreibzeile                 Strut für Tabellenzellen zum Ausfüllen (8 mm hoch)
%  \alt{9}, \neu{7}              überschriebener alter Wert / neuer Wert (Dijkstra-Tabellen)
\usepackage{colortbl,xspace}
\usepackage[normalem]{ulem}
\newcommand{\lsgrest}[1]{\leavevmode\unskip\hspace{1.5mm}\ifloesung\rlap{\raisebox{0.8pt}{\lsgfarbe\,#1}}\fi\leaders\hbox{\rule[-0.6ex]{0.5pt}{0.4pt}}\hfill\null\par}
\newcommand{\restzeile}{\lsgrest{}}
\newcommand{\lsgtext}[1]{\ifloesung\leavevmode{\lsgfarbe#1}\fi}
\newcommand{\lueckentext}{\linespread{1.5}\selectfont\raggedright}
\newcommand{\schreibzeile}{\rule[-2.8mm]{0pt}{8mm}}
\newcommand{\alt}[1]{\leavevmode{\color{blue!60!black}\sout{#1}}}
\newcommand{\neu}[1]{\leavevmode{\color{green!50!black}\bfseries #1}}
\colorlet{djaktuell}{orange!40}
\colorlet{djfertig}{blue!22}

% ---------- Graphen zeichnen ----------
\tikzset{
  knoten/.style={draw,rounded corners=3pt,fill=orange!12,minimum height=6mm,inner sep=3pt,font=\small},
  kknoten/.style={draw,circle,fill=orange!12,minimum size=7mm,inner sep=1pt,font=\small},
  kante/.style={thick},
  gkante/.style={thick,-{Stealth[length=2.6mm]}},
  dkante/.style={thick,{Stealth[length=2.6mm]}-{Stealth[length=2.6mm]}},
  gewicht/.style={font=\scriptsize,fill=white,inner sep=1pt,text=teal!60!black},
  lsgkante/.style={thick,red!75!black},
  lsggkante/.style={thick,red!75!black,-{Stealth[length=2.6mm]}},
  lsgknoten/.style={knoten,draw=red!75!black,text=red!75!black,fill=red!4},
}
% Adjazenzmatrix zum Ausfüllen: \adjmx[Lösungseinträge]{Anzahl}{Knotenliste}{Zellbreite}
% Lösungseinträge als Zeile/Spalte/Text, z. B. \adjmx[1/2/X,2/1/X]{3}{A,B,C}{7mm}
\newcommand{\adjmx}[4][]{\begin{tikzpicture}[x=#4,y=#4]
  \foreach \n [count=\i] in {#3} {\node[font=\small\bfseries] at (\i,0) {\n}; \node[font=\small\bfseries,anchor=east] at (0.45,-\i) {\n};}
  \foreach \k in {0,...,#2} {\draw[gray] (\k+0.5,-0.5) -- (\k+0.5,-#2-0.5); \draw[gray] (0.5,-\k-0.5) -- (#2+0.5,-\k-0.5);}
  \ifloesung\ifstrempty{#1}{}{\foreach \z/\s/\t in {#1} {\node[text=red!75!black,font=\small] at (\s,-\z) {\t};}}\fi
\end{tikzpicture}}

% ---------- Java-Listings ----------
\lstset{language=Java,basicstyle=\ttfamily\small,keywordstyle=\bfseries,commentstyle=\color{gray!80!black},stringstyle=\color{black},showstringspaces=false,columns=flexible,keepspaces=true,tabsize=3,frame=single,framerule=0.4pt,rulecolor=\color{black},xleftmargin=2mm,framexleftmargin=1mm,aboveskip=2mm,belowskip=2mm,breaklines=true,escapechar=§,
 literate={ä}{{\"a}}1 {ö}{{\"o}}1 {ü}{{\"u}}1 {Ä}{{\"A}}1 {Ö}{{\"O}}1 {Ü}{{\"U}}1 {ß}{{\ss}}1 {–}{{--}}1 {„}{{\glqq}}1 {“}{{\grqq}}1}
\lstnewenvironment{java}[1][]{\lstset{#1}}{}

%  Übersetzen mit XeLaTeX, LuaLaTeX, pdfLaTeX oder Tectonic.
%
%  Blatt:   xelatex ab-3-1-grundbegriffe.tex
%  Lösung:  xelatex ab-3-1-grundbegriffe-lsg.tex
%           (Datei enthält nur: \def\mitloesung{}\documentclass[11pt,a4paper]{article}
% =====================================================================
%  Gemeinsame Vorlage der Arbeitsblätter "Graphen" – Informatik 11 (NTG)
%  (gleiches Layout wie informatik/10/java/arbeitsblaetter/ab-vorlage.tex)
%  Einbinden in jedem Blatt direkt nach \documentclass: \input{ab-vorlage}
%  Übersetzen mit XeLaTeX, LuaLaTeX, pdfLaTeX oder Tectonic.
%
%  Blatt:   xelatex ab-3-1-grundbegriffe.tex
%  Lösung:  xelatex ab-3-1-grundbegriffe-lsg.tex
%           (Datei enthält nur: \def\mitloesung{}\input{ab-3-1-grundbegriffe})
%           (füllt alle Lücken rot aus; siehe \lsg, \lsglinien, \lsgoder)
% =====================================================================
\usepackage[a4paper,left=18mm,right=18mm,top=26mm,bottom=17mm,headheight=15mm,headsep=4mm,footskip=8mm]{geometry}
\usepackage{iftex}
\ifPDFTeX
  \usepackage[T1]{fontenc}
  \usepackage[utf8]{inputenc}
  \usepackage{helvet}
  \usepackage[scaled=0.86]{beramono}
  \renewcommand{\familydefault}{\sfdefault}
\else
  \usepackage{fontspec}
  \setmainfont{texgyreheros}[Extension=.otf,UprightFont=*-regular,BoldFont=*-bold,ItalicFont=*-italic,BoldItalicFont=*-bolditalic]
  \setmonofont{DejaVuSansMono}[Extension=.ttf,UprightFont=*,BoldFont=*-Bold,ItalicFont=*-Oblique,BoldItalicFont=*-BoldOblique,Scale=0.86]
\fi
\usepackage[ngerman,shorthands=off]{babel}
\usepackage{xcolor,graphicx,tabularx,array,enumitem,fancyhdr,tikz,listings,multicol,calc,etoolbox}
\usepackage[most]{tcolorbox}
\usetikzlibrary{arrows.meta,positioning,calc,decorations.pathreplacing,shapes.geometric}
\usepackage[hidelinks]{hyperref}
\setlength{\parindent}{0pt}
\setlength{\parskip}{3pt}
\setlist{nosep,leftmargin=*}

% ---------- Lösungsschalter ----------
\newif\ifloesung
\ifdefined\mitloesung\loesungtrue\fi
\newcommand{\lsgfarbe}{\color{red!75!black}}

% ---------- Kopf- und Fußzeile (einheitliche Form) ----------
\newcommand{\lehrkraft}{Hau}
\newcommand{\themenbereich}{3. Graphen}
\newcommand{\abnummer}{}
\pagestyle{fancy}\fancyhf{}
\renewcommand{\headrulewidth}{0.4pt}
\fancyhead[L]{\small Klasse 11\\Datum:}
\fancyhead[C]{\small Themenbereich:\\\themenbereich}
\fancyhead[R]{\small \lehrkraft\ifloesung\\{\lsgfarbe\bfseries Lösung}\fi}
\fancyfoot[C]{\scriptsize edu-mrh.de\,·\,Informatik 11\,·\,Blatt \abnummer\ifloesung\ (Lösung)\fi\,·\,Seite \thepage}
\newcommand{\abtitel}[2]{\renewcommand{\abnummer}{#1}\begin{center}\Large\bfseries #1\quad\underline{#2}\end{center}\vspace{-1mm}}

% ---------- Boxen ----------
\newenvironment{aufgabe}[2][]{\begin{tcolorbox}[enhanced,breakable,colback=white,colframe=black,boxrule=0.7pt,arc=3mm,left=2mm,right=2mm,top=1.5mm,bottom=1.5mm,before skip=2.5mm,after skip=2.5mm]\textbf{Aufgabe #2}\ifstrempty{#1}{}{ \textit{#1}}\par\vspace{0.6mm}}{\end{tcolorbox}}
\newtcolorbox{merke}[1][Merke]{enhanced,breakable,colback=green!5,colframe=green!40!black,boxrule=0.8pt,arc=2mm,left=2mm,right=2mm,top=1.5mm,bottom=1.5mm,before skip=2.5mm,after skip=2.5mm,title={#1},fonttitle=\bfseries,coltitle=white}
\newtcolorbox{hinweis}{enhanced,breakable,colback=yellow!12,colframe=orange!70!black,boxrule=0.6pt,arc=2mm,left=2mm,right=2mm,top=1mm,bottom=1mm,before skip=2mm,after skip=2mm}
\newtcolorbox{ausblick}[1][Ausblick]{enhanced,breakable,colback=teal!5,colframe=teal!60!black,boxrule=0.8pt,arc=2mm,left=2mm,right=2mm,top=1.5mm,bottom=1.5mm,before skip=2.5mm,after skip=2.5mm,title={#1},fonttitle=\bfseries,coltitle=white}
\newcommand{\web}[1]{{\small\color{gray}\textit{Website: #1}}}

% ---------- Lücken, Linien, Karos ----------
\newcommand{\luecke}[1][3.5cm]{\rule[-0.5ex]{#1}{0.4pt}}
\newcommand{\linien}[1]{\par\vspace{1.5mm}\foreach \i in {1,...,#1}{\noindent\rule[0pt]{\linewidth}{0.3pt}\par\vspace{5.5mm}}\vspace{-3mm}}
\newcommand{\kariert}[1]{\par\vspace{1.5mm}\noindent\begin{tikzpicture}\draw[step=5mm,gray!55,thin] (0,0) grid (\linewidth-1pt,#1*5mm);\end{tikzpicture}\par}
\newcommand{\karobox}[2]{\begin{tikzpicture}\draw[step=5mm,gray!55,thin] (0,0) grid (#1,#2);\end{tikzpicture}}
\newcommand{\code}[1]{\texttt{#1}}

% Lücke mit Lösung: \lsg[Mindestbreite]{Lösungstext}
\newlength{\lsgbreite}
\newcommand{\lsg}[2][3.5cm]{\ifloesung\settowidth{\lsgbreite}{\,#2\,}\ifdim\lsgbreite<#1\setlength{\lsgbreite}{#1}\fi\underline{\makebox[\lsgbreite][l]{\lsgfarbe\,#2}}\else\luecke[#1]\fi}
% Schreiblinien mit Lösung: \lsglinien{Anzahl}{Lösungstext} (gleiche Höhe wie \linien)
\newcommand{\lsglinien}[2]{\ifloesung\par\vspace{1.5mm}\noindent\begin{minipage}[t][#1\dimexpr 6.9mm\relax][t]{\linewidth}\lsgfarbe #2\end{minipage}\par\vspace{-1.5mm}\else\linien{#1}\fi}
% Beliebiger Inhalt: \lsgoder{nur in der Lösung}{nur im Blatt}
\newcommand{\lsgoder}[2]{\ifloesung\leavevmode#1\else#2\fi}

% ---------- Ergänzung (aus informatik/10/java/arbeitsblaetter/ab-vorlage.tex) ----------
% Bestehende Makros oben bleiben unverändert, damit ältere Blätter gleich aussehen.
%  \restzeile / \lsgrest{Text}   Lücke bis zum Zeilenende
%  \lsgtext{Text}                nur in der Lösung sichtbar (z. B. Tabellenzellen)
%  \lueckentext                  größerer Zeilenabstand, Flattersatz für Absätze mit Lücken
%  \schreibzeile                 Strut für Tabellenzellen zum Ausfüllen (8 mm hoch)
%  \alt{9}, \neu{7}              überschriebener alter Wert / neuer Wert (Dijkstra-Tabellen)
\usepackage{colortbl,xspace}
\usepackage[normalem]{ulem}
\newcommand{\lsgrest}[1]{\leavevmode\unskip\hspace{1.5mm}\ifloesung\rlap{\raisebox{0.8pt}{\lsgfarbe\,#1}}\fi\leaders\hbox{\rule[-0.6ex]{0.5pt}{0.4pt}}\hfill\null\par}
\newcommand{\restzeile}{\lsgrest{}}
\newcommand{\lsgtext}[1]{\ifloesung\leavevmode{\lsgfarbe#1}\fi}
\newcommand{\lueckentext}{\linespread{1.5}\selectfont\raggedright}
\newcommand{\schreibzeile}{\rule[-2.8mm]{0pt}{8mm}}
\newcommand{\alt}[1]{\leavevmode{\color{blue!60!black}\sout{#1}}}
\newcommand{\neu}[1]{\leavevmode{\color{green!50!black}\bfseries #1}}
\colorlet{djaktuell}{orange!40}
\colorlet{djfertig}{blue!22}

% ---------- Graphen zeichnen ----------
\tikzset{
  knoten/.style={draw,rounded corners=3pt,fill=orange!12,minimum height=6mm,inner sep=3pt,font=\small},
  kknoten/.style={draw,circle,fill=orange!12,minimum size=7mm,inner sep=1pt,font=\small},
  kante/.style={thick},
  gkante/.style={thick,-{Stealth[length=2.6mm]}},
  dkante/.style={thick,{Stealth[length=2.6mm]}-{Stealth[length=2.6mm]}},
  gewicht/.style={font=\scriptsize,fill=white,inner sep=1pt,text=teal!60!black},
  lsgkante/.style={thick,red!75!black},
  lsggkante/.style={thick,red!75!black,-{Stealth[length=2.6mm]}},
  lsgknoten/.style={knoten,draw=red!75!black,text=red!75!black,fill=red!4},
}
% Adjazenzmatrix zum Ausfüllen: \adjmx[Lösungseinträge]{Anzahl}{Knotenliste}{Zellbreite}
% Lösungseinträge als Zeile/Spalte/Text, z. B. \adjmx[1/2/X,2/1/X]{3}{A,B,C}{7mm}
\newcommand{\adjmx}[4][]{\begin{tikzpicture}[x=#4,y=#4]
  \foreach \n [count=\i] in {#3} {\node[font=\small\bfseries] at (\i,0) {\n}; \node[font=\small\bfseries,anchor=east] at (0.45,-\i) {\n};}
  \foreach \k in {0,...,#2} {\draw[gray] (\k+0.5,-0.5) -- (\k+0.5,-#2-0.5); \draw[gray] (0.5,-\k-0.5) -- (#2+0.5,-\k-0.5);}
  \ifloesung\ifstrempty{#1}{}{\foreach \z/\s/\t in {#1} {\node[text=red!75!black,font=\small] at (\s,-\z) {\t};}}\fi
\end{tikzpicture}}

% ---------- Java-Listings ----------
\lstset{language=Java,basicstyle=\ttfamily\small,keywordstyle=\bfseries,commentstyle=\color{gray!80!black},stringstyle=\color{black},showstringspaces=false,columns=flexible,keepspaces=true,tabsize=3,frame=single,framerule=0.4pt,rulecolor=\color{black},xleftmargin=2mm,framexleftmargin=1mm,aboveskip=2mm,belowskip=2mm,breaklines=true,escapechar=§,
 literate={ä}{{\"a}}1 {ö}{{\"o}}1 {ü}{{\"u}}1 {Ä}{{\"A}}1 {Ö}{{\"O}}1 {Ü}{{\"U}}1 {ß}{{\ss}}1 {–}{{--}}1 {„}{{\glqq}}1 {“}{{\grqq}}1}
\lstnewenvironment{java}[1][]{\lstset{#1}}{}

% Dörfer-Graph (Merkkasten Begriffe): Koordinaten wie im Begriffe-Explorer der Website
\newcommand{\doerferknoten}{%
  \node[knoten,font=\tiny,minimum height=4mm,inner sep=1.5pt] (A) at (4.4,3.3) {Altdorf};
  \node[knoten,font=\tiny,minimum height=4mm,inner sep=1.5pt] (W) at (1.2,2.9) {Weiler};
  \node[knoten,font=\tiny,minimum height=4mm,inner sep=1.5pt] (F) at (5.1,2.0) {Fischbach};
  \node[knoten,font=\tiny,minimum height=4mm,inner sep=1.5pt] (Z) at (2.9,2.0) {Ziegelstein};
  \node[knoten,font=\tiny,minimum height=4mm,inner sep=1.5pt] (N) at (0.9,0.95) {Neustadt};
  \node[knoten,font=\tiny,minimum height=4mm,inner sep=1.5pt] (B) at (4.4,0.8) {Burg};
  \node[knoten,font=\tiny,minimum height=4mm,inner sep=1.5pt] (R) at (2.6,0.2) {Rain};}
\begin{document}
\abtitel{3.1}{Graphen -- Grundbegriffe}

Hänsel und Gretel haben die Hexe überlistet. Doch die Brotkrumen, mit denen sie ihren Weg markiert hatten, haben die Vögel gefressen. Findest du trotzdem nach Hause?

\begin{aufgabe}[Hänsel und Gretel]{1}
Starte das Spiel auf der Website (Station 3.1).
\textbf{a)} Zeichne beim Spielen eine Karte des Waldes: jeder Ort ein Kästchen, jeder Weg eine Linie. Führt ein Weg nur in eine Richtung, zeichne einen Pfeil. Das Hexenhaus ist schon eingezeichnet.
\end{aufgabe}
\noindent\begin{tikzpicture}
\draw[step=5mm,gray!55,thin] (0,0) grid (17.4,7);
\node[knoten,fill=green!15] (hex) at (4,2.3) {Hexenhaus (Start)};
\ifloesung
  \node[lsgknoten] (warn) at (5.5,6.5) {Warnung};
  \node[lsgknoten] (wolf) at (10.5,6.5) {Werwolf};
  \node[lsgknoten] (stadt) at (1.6,5.2) {Stadt};
  \node[lsgknoten] (abz) at (5.5,5.2) {Abzweigung};
  \node[lsgknoten] (spinne) at (13,5.2) {Spinne};
  \node[lsgknoten] (baeck) at (1.6,3.6) {Bäckerei};
  \node[lsgknoten] (pilze) at (4,3.9) {Pilze};
  \node[lsgknoten] (w3) at (13,3.4) {Wald3};
  \node[lsgknoten] (w1) at (8,1.0) {Wald1};
  \node[lsgknoten] (stein) at (12.2,2.0) {Steinkreis};
  \node[lsgknoten] (w2) at (15,0.8) {Wald2};
  \node[lsgknoten] (heim) at (16.2,2.6) {Zuhause};
  \draw[lsgkante] (stadt) -- (baeck);
  \draw[lsgkante] (stadt) -- (pilze);
  \draw[lsgkante] (pilze) -- (abz);
  \draw[lsgkante] (pilze) -- (hex);
  \draw[lsggkante] (pilze) to[bend left=50] node[right,font=\scriptsize,text=red!75!black] {Pilze essen} (hex);
  \draw[lsgkante] (abz) -- (warn);
  \draw[lsggkante] (warn) -- (wolf);
  \draw[lsgkante] (abz) -- (spinne);
  \draw[lsgkante] (spinne) -- (w3);
  \draw[lsgkante] (hex) -- (w1);
  \draw[lsgkante] (w1) -- (w3);
  \draw[lsgkante] (w1) -- (w2);
  \draw[lsgkante] (w1) -- (stein);
  \draw[lsgkante] (w2) -- (w3);
  \draw[lsgkante] (w2) -- (stein);
  \draw[lsgkante] (w3) -- (stein);
  \draw[lsggkante] (w2) -- (heim);
  \node[text=red!75!black,font=\scriptsize,align=left,anchor=north west] at (0.2,1.3) {13 Knoten, 17 Kanten\\(14 ungerichtet, 3 gerichtet)};
\fi
\end{tikzpicture}

\textbf{b)} Am Ende zeigt dir das Spiel deinen Weg. Notiere ihn und die Anzahl der Schritte:
\lsglinien{1}{individuell; kürzester Weg: Hexenhaus -- Wald1 -- Wald2 -- Zuhause (3 Schritte)}

\begin{minipage}[t]{0.6\linewidth}
\begin{merke}[Graph und Kurzschreibweise $G = (V, E)$]\raggedright
Ein Graph besteht aus \lsg[2.3cm]{Knoten} (engl. \textit{vertices}) und \lsg[2.3cm]{Kanten} (engl. \textit{edges}). Eine Kante verbindet zwei Knoten. Im Spiel sind die Orte die \lsg[2cm]{Knoten}, die Wege die \lsg[2cm]{Kanten}.\\[1mm]
Kurzschreibweise $G = (V, E)$: $V$ ist die Menge der Knoten, $E$ die Menge der Kanten.
\begin{itemize}
\item $(A, B)$: \lsg[2.5cm]{gerichtete} Kante, nur von $A$ nach $B$
\item $\{A, B\}$: \lsg[2.5cm]{ungerichtete} Kante, steht für $(A, B)$ und \lsg[1.4cm]{$(B, A)$}
\end{itemize}
\end{merke}
\end{minipage}\hfill
\begin{minipage}[t]{0.37\linewidth}
\vspace{1mm}
\centering\small \textbf{Beispiel: Wer folgt wem?}\\[1mm]
\begin{tikzpicture}[x=0.95cm,y=0.9cm]
\node[knoten] (leon) at (0.7,2.2) {Leon};
\node[knoten] (letifa) at (3.3,2.4) {Letifa};
\node[knoten] (luise) at (4.3,1.1) {Luise};
\node[knoten] (hamza) at (2.4,0.3) {Hamza};
\node[knoten] (sabine) at (0.5,0.6) {Sabine};
\draw[dkante] (leon) -- (letifa);
\draw[gkante] (letifa) -- (luise);
\draw[gkante] (leon) -- (luise);
\draw[dkante] (leon) -- (hamza);
\draw[gkante] (hamza) -- (luise);
\end{tikzpicture}
\end{minipage}\\[1mm]
{\small $V = \{\text{Hamza}, \text{Leon}, \text{Letifa}, \text{Luise}, \text{Sabine}\}$\\
$E = \{\{\text{Leon}, \text{Letifa}\}, \{\text{Leon}, \text{Hamza}\}, (\text{Letifa}, \text{Luise}), (\text{Leon}, \text{Luise}), (\text{Hamza}, \text{Luise})\}$}

\begin{aufgabe}[Hänsel und Gretel]{1}
\textbf{c)} Schreibe den Ausschnitt deiner Karte mit den Knoten Hexenhaus, Pilze, Stadt, Bäckerei und Wald1 in Kurzschreibweise auf (5 Kanten).\\[2mm]
$V = \{$\lsg[15.4cm]{Hexenhaus, Pilze, Stadt, Bäckerei, Wald1}$\}$\\[3mm]
$E = \{$\lsg[15.4cm]{\{Hexenhaus, Pilze\}, (Pilze, Hexenhaus), \{Pilze, Stadt\}, \{Stadt, Bäckerei\}, \{Hexenhaus, Wald1\}}$\}$
\end{aufgabe}

\newpage
\begin{minipage}[t]{0.71\linewidth}
\begin{merke}[Pfade, Zyklen, Richtung, Gewicht]\raggedright\small
\textbf{Pfad (Weg):} Folge von Knoten, in der aufeinanderfolgende Knoten durch eine \lsg[1.8cm]{Kante} verbunden sind; keine Kante wird doppelt benutzt.\\[0.6mm]
\textbf{Einfacher Pfad:} Pfad, in dem sich kein \lsg[1.8cm]{Knoten} wiederholt.\\
Beispiel in A: \lsg[5.6cm]{Altdorf -- Weiler -- Ziegelstein -- Neustadt}\\[0.6mm]
\textbf{Länge eines Pfades:} ungewichtet die \lsg[3.2cm]{Anzahl der Kanten}, gewichtet die \lsg[3.2cm]{Summe der Gewichte}.\\[0.6mm]
\textbf{Zyklus:} einfacher Pfad, dessen \lsg[5cm]{Start- und Endknoten gleich sind}.\\
Beispiel in A: \lsg[5.6cm]{Weiler -- Ziegelstein -- Fischbach -- Weiler}\\
Graph mit Zyklus: \emph{zyklisch}, sonst \lsg[3.4cm]{azyklisch (zyklenfrei)}.\\[0.6mm]
\textbf{Gerichtet:} Kanten nur \lsg[3cm]{in Pfeilrichtung} (Pfeil).\ \textbf{Ungerichtet:} \lsg[2.8cm]{in beide Richtungen} (Strich).\\[0.6mm]
\textbf{Stark zusammenhängend:} Von \lsg[2.4cm]{jedem Knoten} gibt es einen Pfad zu \lsg[3.2cm]{jedem anderen Knoten}.\\
\textbf{Schwach zusammenhängend:} \lsg[5.4cm]{nur ohne Beachtung der Richtungen}\\ \lsg[\linewidth]{hängt alles zusammen. B ist schwach, denn von Rain führt keine Kante weg.}\\[0.6mm]
\textbf{Gewichtet:} Jede Kante trägt eine \lsg[2.6cm]{Zahl (Gewicht)}, z.\,B. \lsg[3.6cm]{Entfernung, Fahrzeit}.
\end{merke}
\end{minipage}\hfill
\begin{minipage}[t]{0.27\linewidth}
\vspace{1mm}
\textbf{\small A} {\scriptsize ungerichtet}\\
\begin{tikzpicture}[x=0.8cm,y=0.85cm]
\doerferknoten
\draw[kante] (A)--(W); \draw[kante] (F)--(A); \draw[kante] (W)--(F); \draw[kante] (W)--(Z);
\draw[kante] (Z)--(F); \draw[kante] (Z)--(N); \draw[kante] (B)--(Z); \draw[kante] (F)--(B);
\draw[kante] (N)--(R); \draw[kante] (B)--(R);
\end{tikzpicture}\\[2mm]
\textbf{\small B} {\scriptsize gerichtet}\\
\begin{tikzpicture}[x=0.8cm,y=0.85cm]
\doerferknoten
\draw[gkante] (A)--(W); \draw[gkante] (F)--(A); \draw[gkante] (W)--(F); \draw[gkante] (W)--(Z);
\draw[gkante] (Z)--(F); \draw[gkante] (Z)--(N); \draw[gkante] (B)--(Z); \draw[gkante] (F)--(B);
\draw[gkante] (N)--(R); \draw[gkante] (B)--(R);
\end{tikzpicture}
\end{minipage}

\begin{aufgabe}{2}
{\raggedright Zeichne $G = (V, E)$ mit $V = \{A, B, C, D, E, F\}$ und\\ $E = \{(A, B), (A, C), (A, E), (B, C), (C, C), (C, D), (E, C), (F, D)\}$. Beantworte die Fragen.\par}\vspace{1mm}
\begin{minipage}[t]{0.42\linewidth}
\vspace{0pt}
\begin{tikzpicture}
\draw[step=5mm,gray!55,thin] (0,0) grid (7,3);
\ifloesung
  \foreach \n/\x/\y/\t in {a/0.6/1.5/A, b/2/2.55/B, e/2/0.45/E, c/3.4/1.5/C, d/4.9/1.5/D, f/6.4/1.5/F} {
    \node[kknoten,minimum size=6mm,draw=red!75!black,text=red!75!black] (\n) at (\x,\y) {\t};}
  \draw[lsggkante] (a)--(b); \draw[lsggkante] (a)--(c); \draw[lsggkante] (a)--(e);
  \draw[lsggkante] (b)--(c); \draw[lsggkante] (c)--(d); \draw[lsggkante] (e)--(c);
  \draw[lsggkante] (f)--(d);
  \draw[lsggkante] (c) to[out=60,in=120,looseness=6] (c);
\fi
\end{tikzpicture}
\end{minipage}\hfill
\begin{minipage}[t]{0.55\linewidth}
\vspace{0pt}\small\raggedright
Gerichtet? Woran erkennbar? \lsg[3.9cm]{ja, Paare ( ) statt \{ \}}\\[1.5mm]
Ein Zyklus: \lsg[5.7cm]{(C, C), eine Schlinge (Länge 1)}\\[1.5mm]
Stark/schwach zusammenhängend? \lsg[3.6cm]{schwach: von D}\\[1.5mm]
\lsg[0.98\linewidth]{kommt man nirgends hin.}\\[1.5mm]
Länge des längsten einfachen Pfades: \lsg[3.2cm]{3 (A--B--C--D)}
\end{minipage}
\end{aufgabe}

\begin{aufgabe}[Zurück in den Wald]{3}
Deine Karte aus Aufgabe 1: \textbf{a)} gerichtet?\ \textbf{b)} ein Zyklus?\ \textbf{c)} stark oder schwach zusammenhängend?\ \textbf{d)} „Der Weg zurück … scheint verschwunden“: Ist die Kante Hexenhaus -- Wald1 gerichtet? Prüfe im Spiel.
\lsglinien{2}{\small a) gemischt; gerichtet nur (Pilze, Hexenhaus), (Warnung, Werwolf), (Wald2, Zuhause)\quad b) Wald1 -- Wald2 -- Wald3 -- Wald1\\ c) schwach: Von Zuhause/Werwolf führt keine Kante weg.\quad d) nein: Von Wald2 aus führt ein Weg zurück.}
\end{aufgabe}

\begin{aufgabe}[Ordne die Graphen ein.]{4}
\vspace{-1mm}
\small\renewcommand{\arraystretch}{1.15}
\begin{tabularx}{\linewidth}{|>{\centering\arraybackslash}m{5.6cm}|>{\centering\arraybackslash}m{1.7cm}|>{\centering\arraybackslash}m{1.1cm}|>{\centering\arraybackslash}m{1.1cm}|>{\centering\arraybackslash}m{1.1cm}|>{\centering\arraybackslash}X|}\hline
\textbf{Graph} & \scriptsize\textbf{zusammen-\newline hängend?}\newline stark/schwach/nein & \scriptsize\textbf{ge-\newline richtet?} & \scriptsize\textbf{ge-\newline wichtet?} & \scriptsize\textbf{zyklisch?} & \scriptsize\textbf{Länge eines längsten einfachen Pfades}\\\hline
\rule{0pt}{8mm}\textbf{I}\ \begin{tikzpicture}[baseline=(k5.base),x=0.85cm,y=0.55cm]
\foreach \n/\x/\y in {1/0/1,2/1/1,3/2/1,4/3/1,5/1/0,6/3/0} {\node[kknoten,minimum size=5mm,inner sep=0pt,font=\tiny] (k\n) at (\x,\y) {K\n};}
\draw[kante] (k1)--(k2)--(k3)--(k4); \draw[kante] (k2)--(k5);
\end{tikzpicture}
& \lsgoder{\lsgfarbe nein}{} & \lsgoder{\lsgfarbe nein}{} & \lsgoder{\lsgfarbe nein}{} & \lsgoder{\lsgfarbe nein}{} & \lsgoder{\lsgfarbe 3: K1--K2--K3--K4}{}\\\hline
\rule{0pt}{5mm}\textbf{II}\ $V = \{A, B, C\}$\newline $E =\{(A, B), (B, A), (C, B), (C, C)\}$
& \lsgoder{\lsgfarbe schwach}{} & \lsgoder{\lsgfarbe ja}{} & \lsgoder{\lsgfarbe nein}{} & \lsgoder{\lsgfarbe ja}{} & \lsgoder{\lsgfarbe 2: C$\to$B$\to$A}{}\\\hline
\rule{0pt}{8mm}\textbf{III}\ \begin{tikzpicture}[baseline=(k3.base),x=0.9cm,y=0.6cm]
\foreach \n/\x/\y in {1/0/1,2/2/1,3/1/0,4/3/0} {\node[kknoten,minimum size=5mm,inner sep=0pt,font=\tiny] (k\n) at (\x,\y) {K\n};}
\draw[kante] (k1) -- node[gewicht] {4} (k2); \draw[kante] (k2) -- node[gewicht] {1} (k3);
\draw[kante] (k3) -- node[gewicht] {5} (k1); \draw[kante] (k3) -- node[gewicht] {2} (k4);
\end{tikzpicture}
& \lsgoder{\lsgfarbe stark}{} & \lsgoder{\lsgfarbe nein}{} & \lsgoder{\lsgfarbe ja}{} & \lsgoder{\lsgfarbe ja}{} & \lsgoder{\lsgfarbe 11: K2--K1--K3--K4}{}\\\hline
\end{tabularx}
\end{aufgabe}
\end{document}
)
%           (füllt alle Lücken rot aus; siehe \lsg, \lsglinien, \lsgoder)
% =====================================================================
\usepackage[a4paper,left=18mm,right=18mm,top=26mm,bottom=17mm,headheight=15mm,headsep=4mm,footskip=8mm]{geometry}
\usepackage{iftex}
\ifPDFTeX
  \usepackage[T1]{fontenc}
  \usepackage[utf8]{inputenc}
  \usepackage{helvet}
  \usepackage[scaled=0.86]{beramono}
  \renewcommand{\familydefault}{\sfdefault}
\else
  \usepackage{fontspec}
  \setmainfont{texgyreheros}[Extension=.otf,UprightFont=*-regular,BoldFont=*-bold,ItalicFont=*-italic,BoldItalicFont=*-bolditalic]
  \setmonofont{DejaVuSansMono}[Extension=.ttf,UprightFont=*,BoldFont=*-Bold,ItalicFont=*-Oblique,BoldItalicFont=*-BoldOblique,Scale=0.86]
\fi
\usepackage[ngerman,shorthands=off]{babel}
\usepackage{xcolor,graphicx,tabularx,array,enumitem,fancyhdr,tikz,listings,multicol,calc,etoolbox}
\usepackage[most]{tcolorbox}
\usetikzlibrary{arrows.meta,positioning,calc,decorations.pathreplacing,shapes.geometric}
\usepackage[hidelinks]{hyperref}
\setlength{\parindent}{0pt}
\setlength{\parskip}{3pt}
\setlist{nosep,leftmargin=*}

% ---------- Lösungsschalter ----------
\newif\ifloesung
\ifdefined\mitloesung\loesungtrue\fi
\newcommand{\lsgfarbe}{\color{red!75!black}}

% ---------- Kopf- und Fußzeile (einheitliche Form) ----------
\newcommand{\lehrkraft}{Hau}
\newcommand{\themenbereich}{3. Graphen}
\newcommand{\abnummer}{}
\pagestyle{fancy}\fancyhf{}
\renewcommand{\headrulewidth}{0.4pt}
\fancyhead[L]{\small Klasse 11\\Datum:}
\fancyhead[C]{\small Themenbereich:\\\themenbereich}
\fancyhead[R]{\small \lehrkraft\ifloesung\\{\lsgfarbe\bfseries Lösung}\fi}
\fancyfoot[C]{\scriptsize edu-mrh.de\,·\,Informatik 11\,·\,Blatt \abnummer\ifloesung\ (Lösung)\fi\,·\,Seite \thepage}
\newcommand{\abtitel}[2]{\renewcommand{\abnummer}{#1}\begin{center}\Large\bfseries #1\quad\underline{#2}\end{center}\vspace{-1mm}}

% ---------- Boxen ----------
\newenvironment{aufgabe}[2][]{\begin{tcolorbox}[enhanced,breakable,colback=white,colframe=black,boxrule=0.7pt,arc=3mm,left=2mm,right=2mm,top=1.5mm,bottom=1.5mm,before skip=2.5mm,after skip=2.5mm]\textbf{Aufgabe #2}\ifstrempty{#1}{}{ \textit{#1}}\par\vspace{0.6mm}}{\end{tcolorbox}}
\newtcolorbox{merke}[1][Merke]{enhanced,breakable,colback=green!5,colframe=green!40!black,boxrule=0.8pt,arc=2mm,left=2mm,right=2mm,top=1.5mm,bottom=1.5mm,before skip=2.5mm,after skip=2.5mm,title={#1},fonttitle=\bfseries,coltitle=white}
\newtcolorbox{hinweis}{enhanced,breakable,colback=yellow!12,colframe=orange!70!black,boxrule=0.6pt,arc=2mm,left=2mm,right=2mm,top=1mm,bottom=1mm,before skip=2mm,after skip=2mm}
\newtcolorbox{ausblick}[1][Ausblick]{enhanced,breakable,colback=teal!5,colframe=teal!60!black,boxrule=0.8pt,arc=2mm,left=2mm,right=2mm,top=1.5mm,bottom=1.5mm,before skip=2.5mm,after skip=2.5mm,title={#1},fonttitle=\bfseries,coltitle=white}
\newcommand{\web}[1]{{\small\color{gray}\textit{Website: #1}}}

% ---------- Lücken, Linien, Karos ----------
\newcommand{\luecke}[1][3.5cm]{\rule[-0.5ex]{#1}{0.4pt}}
\newcommand{\linien}[1]{\par\vspace{1.5mm}\foreach \i in {1,...,#1}{\noindent\rule[0pt]{\linewidth}{0.3pt}\par\vspace{5.5mm}}\vspace{-3mm}}
\newcommand{\kariert}[1]{\par\vspace{1.5mm}\noindent\begin{tikzpicture}\draw[step=5mm,gray!55,thin] (0,0) grid (\linewidth-1pt,#1*5mm);\end{tikzpicture}\par}
\newcommand{\karobox}[2]{\begin{tikzpicture}\draw[step=5mm,gray!55,thin] (0,0) grid (#1,#2);\end{tikzpicture}}
\newcommand{\code}[1]{\texttt{#1}}

% Lücke mit Lösung: \lsg[Mindestbreite]{Lösungstext}
\newlength{\lsgbreite}
\newcommand{\lsg}[2][3.5cm]{\ifloesung\settowidth{\lsgbreite}{\,#2\,}\ifdim\lsgbreite<#1\setlength{\lsgbreite}{#1}\fi\underline{\makebox[\lsgbreite][l]{\lsgfarbe\,#2}}\else\luecke[#1]\fi}
% Schreiblinien mit Lösung: \lsglinien{Anzahl}{Lösungstext} (gleiche Höhe wie \linien)
\newcommand{\lsglinien}[2]{\ifloesung\par\vspace{1.5mm}\noindent\begin{minipage}[t][#1\dimexpr 6.9mm\relax][t]{\linewidth}\lsgfarbe #2\end{minipage}\par\vspace{-1.5mm}\else\linien{#1}\fi}
% Beliebiger Inhalt: \lsgoder{nur in der Lösung}{nur im Blatt}
\newcommand{\lsgoder}[2]{\ifloesung\leavevmode#1\else#2\fi}

% ---------- Ergänzung (aus informatik/10/java/arbeitsblaetter/ab-vorlage.tex) ----------
% Bestehende Makros oben bleiben unverändert, damit ältere Blätter gleich aussehen.
%  \restzeile / \lsgrest{Text}   Lücke bis zum Zeilenende
%  \lsgtext{Text}                nur in der Lösung sichtbar (z. B. Tabellenzellen)
%  \lueckentext                  größerer Zeilenabstand, Flattersatz für Absätze mit Lücken
%  \schreibzeile                 Strut für Tabellenzellen zum Ausfüllen (8 mm hoch)
%  \alt{9}, \neu{7}              überschriebener alter Wert / neuer Wert (Dijkstra-Tabellen)
\usepackage{colortbl,xspace}
\usepackage[normalem]{ulem}
\newcommand{\lsgrest}[1]{\leavevmode\unskip\hspace{1.5mm}\ifloesung\rlap{\raisebox{0.8pt}{\lsgfarbe\,#1}}\fi\leaders\hbox{\rule[-0.6ex]{0.5pt}{0.4pt}}\hfill\null\par}
\newcommand{\restzeile}{\lsgrest{}}
\newcommand{\lsgtext}[1]{\ifloesung\leavevmode{\lsgfarbe#1}\fi}
\newcommand{\lueckentext}{\linespread{1.5}\selectfont\raggedright}
\newcommand{\schreibzeile}{\rule[-2.8mm]{0pt}{8mm}}
\newcommand{\alt}[1]{\leavevmode{\color{blue!60!black}\sout{#1}}}
\newcommand{\neu}[1]{\leavevmode{\color{green!50!black}\bfseries #1}}
\colorlet{djaktuell}{orange!40}
\colorlet{djfertig}{blue!22}

% ---------- Graphen zeichnen ----------
\tikzset{
  knoten/.style={draw,rounded corners=3pt,fill=orange!12,minimum height=6mm,inner sep=3pt,font=\small},
  kknoten/.style={draw,circle,fill=orange!12,minimum size=7mm,inner sep=1pt,font=\small},
  kante/.style={thick},
  gkante/.style={thick,-{Stealth[length=2.6mm]}},
  dkante/.style={thick,{Stealth[length=2.6mm]}-{Stealth[length=2.6mm]}},
  gewicht/.style={font=\scriptsize,fill=white,inner sep=1pt,text=teal!60!black},
  lsgkante/.style={thick,red!75!black},
  lsggkante/.style={thick,red!75!black,-{Stealth[length=2.6mm]}},
  lsgknoten/.style={knoten,draw=red!75!black,text=red!75!black,fill=red!4},
}
% Adjazenzmatrix zum Ausfüllen: \adjmx[Lösungseinträge]{Anzahl}{Knotenliste}{Zellbreite}
% Lösungseinträge als Zeile/Spalte/Text, z. B. \adjmx[1/2/X,2/1/X]{3}{A,B,C}{7mm}
\newcommand{\adjmx}[4][]{\begin{tikzpicture}[x=#4,y=#4]
  \foreach \n [count=\i] in {#3} {\node[font=\small\bfseries] at (\i,0) {\n}; \node[font=\small\bfseries,anchor=east] at (0.45,-\i) {\n};}
  \foreach \k in {0,...,#2} {\draw[gray] (\k+0.5,-0.5) -- (\k+0.5,-#2-0.5); \draw[gray] (0.5,-\k-0.5) -- (#2+0.5,-\k-0.5);}
  \ifloesung\ifstrempty{#1}{}{\foreach \z/\s/\t in {#1} {\node[text=red!75!black,font=\small] at (\s,-\z) {\t};}}\fi
\end{tikzpicture}}

% ---------- Java-Listings ----------
\lstset{language=Java,basicstyle=\ttfamily\small,keywordstyle=\bfseries,commentstyle=\color{gray!80!black},stringstyle=\color{black},showstringspaces=false,columns=flexible,keepspaces=true,tabsize=3,frame=single,framerule=0.4pt,rulecolor=\color{black},xleftmargin=2mm,framexleftmargin=1mm,aboveskip=2mm,belowskip=2mm,breaklines=true,escapechar=§,
 literate={ä}{{\"a}}1 {ö}{{\"o}}1 {ü}{{\"u}}1 {Ä}{{\"A}}1 {Ö}{{\"O}}1 {Ü}{{\"U}}1 {ß}{{\ss}}1 {–}{{--}}1 {„}{{\glqq}}1 {“}{{\grqq}}1}
\lstnewenvironment{java}[1][]{\lstset{#1}}{}


% ---------- Freibad-Graph (Aufgabe 2), klein, ein Bild je Schritt ----------
% \fbgraph{aktueller Knoten}{fertige Knoten}{Kanten, die geprüft werden}{Beschriftungen}
% Beschriftungen mit \lH{…}, \lK{…}, \lW{…}, \lE{…}, \lF{…} (Pfadlänge am Knoten)
\tikzset{
  dj/.style={draw,circle,minimum size=5.6mm,inner sep=0pt,font=\scriptsize\bfseries,fill=white},
  djkante/.style={gray!70!black,line width=0.6pt},
  djpruef/.style={orange!85!black,line width=1.8pt},
  djgew/.style={font=\scriptsize,inner sep=0.6pt,fill=white,text=black!80},
  djl/.style={font=\scriptsize\bfseries,text=blue!60!black,inner sep=1pt},
}
\newcommand{\lH}[1]{\node[djl,anchor=east] at (H.west) {#1};}
\newcommand{\lK}[1]{\node[djl,anchor=north] at (K.south) {#1};}
\newcommand{\lW}[1]{\node[djl,anchor=south] at (W.north) {#1};}
\newcommand{\lE}[1]{\node[djl,anchor=south] at (E.north) {#1};}
\newcommand{\lF}[1]{\node[djl,anchor=north] at (F.south) {#1};}
\newcommand{\fbgraph}[4]{\begin{tikzpicture}[x=1.05cm,y=0.68cm,baseline=(current bounding box.center)]
  \foreach \n/\x/\y in {H/0/0.75, W/1.35/1.55, E/2.75/1.55, K/2.2/0, F/3.85/0.75} {\coordinate (\n) at (\x,\y);}
  \foreach \a/\b/\w/\p in {H/W/9/left, H/K/2/below, W/E/5/above, W/K/5/left, K/E/12/right, E/F/3/above right, K/F/15/below right} {\draw[djkante] (\a) -- node[djgew,\p,pos=0.5] {\w} (\b);}
  \foreach \a/\b in {#3} {\draw[djpruef] (\a) -- (\b);}
  \foreach \n in {H,W,E,K,F} {\node[dj] (\n) at (\n) {\n};}
  \foreach \n in {#2} {\node[dj,fill=djfertig] at (\n) {\n};}
  \ifstrempty{#1}{}{\node[dj,fill=djaktuell,draw=orange!80!black,line width=1pt] at (#1) {#1};}
  #4
\end{tikzpicture}}
% Tabelle eines Schritts: \fbtab{Farben H K W E F}{H}{K}{W}{E}{F}
% Farbe je Knoten: a = aktuell (orange), f = fertig (blau), {} = ohne
\newcommand{\fbfarbe}[1]{\ifstrequal{#1}{a}{\cellcolor{djaktuell}}{\ifstrequal{#1}{f}{\cellcolor{djfertig}}{}}}
\newcommand{\fbkopf}[5]{\fbfarbe{#1}\textbf{H} & \fbfarbe{#2}\textbf{K} & \fbfarbe{#3}\textbf{W} & \fbfarbe{#4}\textbf{E} & \fbfarbe{#5}\textbf{F}}
\newcommand{\fbtab}[6]{{\renewcommand{\arraystretch}{1}\setlength{\tabcolsep}{2pt}%
  \begin{tabularx}{\linewidth}{|*{5}{>{\centering\arraybackslash}X|}}\hline
  \expandafter\fbkopf#1\\\hline
  \schreibzeile #2 & #3 & #4 & #5 & #6\\\hline
  \end{tabularx}}}
% Kopffarben in Schritt 4 und 5: im Blatt färbt man selbst, in der Lösung vorgegeben
\ifloesung\def\fbvier{fffa{}}\def\fbfuenf{ffffa}\else\def\fbvier{fff{}{}}\def\fbfuenf{{}{}{}{}{}}\fi
\newcommand{\schritt}[3]{\par\noindent\begin{minipage}[c]{5cm}\centering #2\end{minipage}\hfill
  \begin{minipage}[c]{\linewidth-5.2cm}\textbf{Schritt #1}\par\vspace{0.6mm}#3\end{minipage}\par\vspace{1.2mm}}

% ---------- Icons für die Karte (Aufgabe 1) ----------
\newcommand{\ihaus}{\begin{tikzpicture}[scale=0.32,orange!85!black,line width=0.6pt]\draw (0,0) rectangle (1.2,0.9); \draw (-0.15,0.85) -- (0.6,1.45) -- (1.35,0.85); \draw (0.45,0) rectangle (0.75,0.45);\end{tikzpicture}}
\newcommand{\iwald}{\begin{tikzpicture}[scale=0.32,green!40!black,line width=0.6pt]\foreach \x/\h in {0/1.1,0.75/1.5,1.5/1.1}{\draw (\x-0.35,0.2) -- (\x,0.2+\h) -- (\x+0.35,0.2) -- cycle; \draw (\x,0) -- (\x,0.2);}\end{tikzpicture}}
\newcommand{\ikirche}{\begin{tikzpicture}[scale=0.32,black!80,line width=0.6pt]\draw (0,0) rectangle (1.2,0.8); \draw (-0.1,0.75) -- (0.6,1.3) -- (1.3,0.75); \draw (0.6,1.3) -- (0.6,1.75) (0.4,1.55) -- (0.8,1.55); \draw (0.45,0) rectangle (0.75,0.45);\end{tikzpicture}}
\newcommand{\ieis}{\begin{tikzpicture}[scale=0.32,brown!70!black,line width=0.6pt]\draw (0.25,0.8) -- (0.6,-0.2) -- (0.95,0.8); \draw (0.4,1.05) circle (0.27); \draw (0.8,1.05) circle (0.27); \draw (0.6,1.4) circle (0.27);\end{tikzpicture}}
\newcommand{\ibad}{\begin{tikzpicture}[scale=0.32,cyan!60!black,line width=0.6pt]\foreach \y in {0,0.35}{\draw (0,\y) sin (0.25,\y+0.15) cos (0.5,\y) sin (0.75,\y+0.15) cos (1,\y) sin (1.25,\y+0.15) cos (1.5,\y);} \draw (1.1,0.7) -- (1.1,1.3) -- (1.5,1.3); \draw (0.2,0.9) -- (0.9,1.3);\end{tikzpicture}}
\tikzset{ort/.style={draw,rounded corners=2pt,minimum width=11mm,minimum height=8.5mm,inner sep=1pt,fill=white},
         ortname/.style={font=\scriptsize,inner sep=1pt}}

% ---------- Städte (Aufgabe 4) ----------
\tikzset{stadt/.style={draw,ellipse,inner sep=1.2pt,font=\scriptsize,fill=white}}

\begin{document}
\abtitel{3.4}{Der Dijkstra-Algorithmus}

\begin{aufgabe}[Ab ins Freibad]{1}
\begin{minipage}[t]{0.55\linewidth}
\vspace{0pt}\raggedright
Felix will mit seinen Freunden ins Freibad gehen. Einige von ihnen sind bereits da.\\[1mm]
\textbf{a)} Bestimme die kürzeste Weglänge vom Haus zum Freibad.\\[3.5mm]
Weg: \lsg[3.9cm]{H – K – W – E – F}\enspace Länge: \lsg[1.5cm]{15}\,km\\[1.5mm]
\textbf{b)} Am Samstag muss Felix zuvor noch in den Wald, um seinem Opa zu helfen. Wie lang ist sein Weg jetzt?\\[3.5mm]
Weg: \lsg[3.9cm]{H – K – W – E – F}\enspace Länge: \lsg[1.5cm]{15}\,km
\end{minipage}\hfill
\begin{minipage}[t]{0.43\linewidth}
\vspace{0pt}\centering
\begin{tikzpicture}[x=1cm,y=1cm]
  \node[ort,draw=orange!85!black] (H) at (0,1.2) {\ihaus};
  \node[ort] (W) at (1.7,2.35) {\iwald};
  \node[ort] (E) at (4.3,2.35) {\ieis};
  \node[ort,draw=cyan!60!black] (F) at (5.6,1.2) {\ibad};
  \node[ort] (K) at (2.9,0) {\ikirche};
  \foreach \a/\b/\w/\p in {H/W/9 km/above left, H/K/2 km/below left, W/E/5 km/above, W/K/5 km/left, K/E/12 km/right, E/F/3 km/above right, K/F/15 km/below right}
    {\draw[djkante] (\a) -- node[djgew,\p,pos=0.42] {\w} (\b);}
  \node[ortname,below=0pt of H] {Haus (H)};
  \node[ortname,above=0pt of W] {Wald (W)};
  \node[ortname,above=0pt of E] {Eisdiele (E)};
  \node[ortname,below=0pt of F] {Freibad (F)};
  \node[ortname,below=0pt of K] {Kirche (K)};
\end{tikzpicture}
\end{minipage}
\end{aufgabe}

\begin{aufgabe}[Schritt für Schritt]{2}
\raggedright
Notiere den Ablauf in den Tabellen. Knoten H ist der Startpunkt, Schritt 1 und 2 sind schon eingetragen. Du kannst in jedem Schritt nur die Knoten bzw. die Länge zu ihnen (ausgehend von H) updaten, wenn es eine \underline{direkte} Kante ausgehend vom aktuell besuchten Knoten gibt (dicke Kanten im Bild). Wird eine Länge kleiner, streiche die alte durch.\\[1mm]
{\small\colorbox{djaktuell}{\strut orange} aktuell besuchter Knoten\qquad
\colorbox{djfertig}{\strut blau} fertiger Knoten\qquad
\alt{9}~\neu{7} alte und neue Länge}
\end{aufgabe}
\vspace{-1mm}
\noindent\hspace*{5.3cm}{\small\textbf{Knoten und Pfadlänge ausgehend von H}}\hfill\web{Schritt für Schritt}\par\vspace{0.5mm}
\schritt{1}{\fbgraph{H}{}{H/W,H/K}{\lH{0}\lK{2}\lW{9}}}
  {\fbtab{a{}{}{}{}}{0}{2}{9}{$\infty$}{$\infty$}}
\schritt{2}{\fbgraph{K}{H}{K/W,K/E,K/F}{\lH{0}\lK{2}\lW{\alt{9}\,\neu{7}}\lE{14}\lF{17}}}
  {\fbtab{fa{}{}{}}{0}{2}{\alt{9}\ \neu{7}}{14}{17}}
\schritt{3}{\fbgraph{W}{H,K}{W/E}{\lH{0}\lK{2}\lW{7}\ifloesung\lE{\lsgfarbe\sout{14}\,12}\lF{\lsgfarbe 17}\fi}}
  {\fbtab{ffa{}{}}{0}{2}{\lsgtext{7}}{\lsgtext{\sout{14}\ 12}}{\lsgtext{17}}}
\schritt{4}{\fbgraph{E}{H,K,W}{E/F}{\lH{0}\lK{2}\lW{7}\ifloesung\lE{\lsgfarbe 12}\lF{\lsgfarbe\sout{17}\,15}\fi}}
  {\fbtab{\fbvier}{\lsgtext{0}}{\lsgtext{2}}{\lsgtext{7}}{\lsgtext{12}}{\lsgtext{\sout{17}\ 15}}}
\schritt{5}{\fbgraph{F}{H,K,W,E}{}{\lH{0}\lK{2}\lW{7}\ifloesung\lE{\lsgfarbe 12}\lF{\lsgfarbe 15}\fi}}
  {\fbtab{\fbfuenf}{\lsgtext{0}}{\lsgtext{2}}{\lsgtext{7}}{\lsgtext{12}}{\lsgtext{15}}}

\newpage
\begin{aufgabe}[Der Algorithmus]{3}
Bringe die Schritte des Dijkstra-Algorithmus in eine sinnvolle Reihenfolge. Die Buchstaben ergeben ein Lösungswort.\\[1mm]
\begin{tikzpicture}[karte/.style={draw,rounded corners=2pt,fill=white,align=left,font=\small,inner sep=3.5pt,anchor=north west},bst/.style={font=\large\bfseries}]
  \path (0,0) (\linewidth-6mm,-4.05);
  \node[karte,text width=3.0cm] at (0,0) {{\bfseries R} Alle anderen Knoten erhalten die Pfadlänge „unendlich“ ($\infty$).};
  \node[karte,text width=2.4cm] at (3.65,-0.3) {{\bfseries W} Wähle einen beliebigen Startknoten.};
  \node[karte,text width=5.7cm] at (6.65,0) {{\bfseries L} Wähle unter den noch nicht fertigen Knoten den mit der kleinsten Pfadlänge (orange). Aktualisiere für jeden seiner nicht fertigen Nachbarknoten die Pfadlänge, wenn sich über ihn eine kleinere ergibt.};
  \node[karte,text width=5.4cm] at (0.9,-2.1) {{\bfseries D} Markiere den Knoten als fertig (blau). Wiederhole die letzten beiden Schritte, bis alle Knoten fertig sind bzw. der gewünschte Knoten fertig ist.};
  \node[karte,text width=2.7cm] at (13.0,-0.45) {{\bfseries O} Der Startknoten erhält die Pfadlänge 0.};
  \node[anchor=north west,font=\small] at (7.4,-3.1) {\texttt{Hello} \lsg[4.5cm]{\texttt{WORLD}}};
\end{tikzpicture}
\end{aufgabe}

\begin{aufgabe}[Paketdienst]{4}
\begin{minipage}[t]{0.55\linewidth}
\vspace{0pt}\raggedright
Ein Transportunternehmen muss seine Pakete möglichst schnell ausliefern.\\[0.5mm]
\textbf{a)} Ermittle mithilfe des Dijkstra-Algorithmus die kürzeste Verbindungsstrecke zwischen Frankfurt und München. Gib auch die Besuchsreihenfolge an.\\[0.5mm]
\textbf{b)} Recherchiere im Internet oder probiere auf der Website aus, ob der Algorithmus bei der jeweiligen Eigenschaft ausführbar ist:
(i) gerichteter Graph, (ii) negative Kantengewichte, (iii) nicht zusammenhängender Graph.
\end{minipage}\hfill
\begin{minipage}[t]{0.43\linewidth}
\vspace{0pt}\centering
\begin{tikzpicture}[x=0.9cm,y=0.78cm]
  \foreach \n/\t/\x/\y in {Fr/Frankfurt/2.5/3.3, Ma/Mannheim/0.4/2.6, Wu/Würzburg/4.6/2.75, Er/Erfurt/6.6/3.9, Ka/Karlsruhe/1.2/1.35, Nu/Nürnberg/6.3/1.45, Au/Augsburg/3.3/0.55, Mu/München/5.7/-0.05}
    {\coordinate (\n) at (\x,\y);}
  \foreach \a/\b/\w/\p in {Fr/Ma/85/above left, Fr/Wu/119/below left, Wu/Er/193/below right, Wu/Nu/109/right, Wu/Au/150/left, Ma/Ka/70/left, Ka/Au/230/below left, Au/Mu/80/below left, Nu/Mu/170/right}
    {\draw[djkante] (\a) -- node[djgew,\p,pos=0.5] {\w\,km} (\b);}
  \foreach \n/\t/\x/\y in {Fr/Frankfurt/2.5/3.3, Ma/Mannheim/0.4/2.6, Wu/Würzburg/4.6/2.75, Er/Erfurt/6.6/3.9, Ka/Karlsruhe/1.2/1.35, Nu/Nürnberg/6.3/1.45, Au/Augsburg/3.3/0.55, Mu/München/5.7/-0.05}
    {\node[stadt] at (\n) {\t};}
\end{tikzpicture}
\end{minipage}\par\vspace{1mm}
\textbf{a)}\hfill{\small\textbf{Knoten und Pfadlänge ausgehend von Frankfurt}}\hfill\null\par\vspace{0.5mm}
{\setlength{\tabcolsep}{1.5pt}\renewcommand{\arraystretch}{1}\small
\newcommand{\z}[8]{\schreibzeile #1 & #2 & #3 & #4 & #5 & #6 & #7 & #8\\\hline}
\newcommand{\A}{\cellcolor{djaktuell}}\newcommand{\F}{\cellcolor{djfertig}}
\newcommand{\Al}{\ifloesung\cellcolor{djaktuell}\fi}\newcommand{\Fl}{\ifloesung\cellcolor{djfertig}\fi}
\newcommand{\W}[1]{\lsgtext{#1}}\newcommand{\I}{\lsgtext{$\infty$}}
\begin{tabularx}{\linewidth}{|*{8}{>{\centering\arraybackslash}X|}}\hline
\textbf{Frankfurt} & \textbf{Mannheim} & \textbf{Würzburg} & \textbf{Karlsruhe} & \textbf{Erfurt} & \textbf{Augsburg} & \textbf{Nürnberg} & \textbf{München}\\\hline
\z{\A\W{0}}{\W{85}}{\W{119}}{\I}{\I}{\I}{\I}{\I}
\z{\F\W{0}}{\A\W{85}}{\W{119}}{\W{155}}{\I}{\I}{\I}{\I}
\z{\Fl\W{0}}{\Fl\W{85}}{\Al\W{119}}{\W{155}}{\W{312}}{\W{269}}{\W{228}}{\I}
\z{\Fl\W{0}}{\Fl\W{85}}{\Fl\W{119}}{\Al\W{155}}{\W{312}}{\W{269}}{\W{228}}{\I}
\z{\Fl\W{0}}{\Fl\W{85}}{\Fl\W{119}}{\Fl\W{155}}{\W{312}}{\W{269}}{\Al\W{228}}{\W{398}}
\z{\Fl\W{0}}{\Fl\W{85}}{\Fl\W{119}}{\Fl\W{155}}{\W{312}}{\Al\W{269}}{\Fl\W{228}}{\W{\sout{398}\ 349}}
\z{\Fl\W{0}}{\Fl\W{85}}{\Fl\W{119}}{\Fl\W{155}}{\Al\W{312}}{\Fl\W{269}}{\Fl\W{228}}{\W{349}}
\z{\Fl\W{0}}{\Fl\W{85}}{\Fl\W{119}}{\Fl\W{155}}{\Fl\W{312}}{\Fl\W{269}}{\Fl\W{228}}{\Al\W{349}}
\end{tabularx}}\\[1.5mm]
{\lueckentext
Besuchsreihenfolge:\lsgrest{\small Frankfurt, Mannheim, Würzburg, Karlsruhe, Nürnberg, Augsburg, Erfurt, München}
Kürzeste Strecke: \lsg[8cm]{Frankfurt – Würzburg – Augsburg – München}\quad Länge: \lsg[2cm]{349}\,km\par
\textbf{b)} (i) gerichtet:\lsgrest{ja, Kanten werden nur in Pfeilrichtung benutzt.}
(ii) negative Gewichte:\lsgrest{nein, ein fertiger Knoten könnte später noch kürzer werden.}
(iii) nicht zusammenhängend:\lsgrest{ja, nicht erreichbare Knoten behalten die Pfadlänge $\infty$.}
}
\end{aufgabe}
\end{document}
